Network slicing lets enhanced mobile broadband (eMBB), ultra-reliable low-latency communication (URLLC) and massive machine-type communication (mMTC) services share one infrastructure. Most slicing controllers only react once demand has already changed. In this thesis we ask two questions: does it help to predict traffic before bandwidth is allocated, and is a decentralised multi-agent controller a better design than one central controller?

We built a pipeline of self-supervised learning (SSL), a temporal graph neural network (TGNN) forecaster and multi-agent reinforcement learning (MARL), and tested it in a Digital Twin. No public dataset combines topology, per-slice traffic, quality of service and allocation actions, so we made a unified slicing environment from real Telecom Italia Milan traffic (62 days, 10{,}000 cells), with latency calibrated on a public 6G RAN telemetry dataset. All results are measured on a held-out week, and the controllers were chosen on separate validation days.

A MARL controller with shared parameters keeps one actor of 19{,}332 parameters from 9 to 100 cells. A single central agent grows to 604{,}048 parameters and its reward falls from about 0.43 to 0.27. We also compared with rule-based controllers whose targets were tuned on the training days. With a long training budget and averaged seeds, the off-policy learned controller beats them on all seven test days at every size we tried, by 0.023 to 0.026 in reward. It only matches a richer rule that is tuned for each time of day and cell class. An on-policy (PPO-family) learner does beat that rule on every test day, by 0.009 in reward at 25 cells.

A TGNN trained on 39 days beats persistence on all 10{,}000 cells (its mean squared error is lower by 18\%, 33\% and 31\% for eMBB, URLLC and mMTC). Giving its forecast to the rule-based controller raises the reward by about 0.03, which closes about one fifth of the gap to perfect knowledge of the next interval. We have not shown that the forecast helps the learned controller. Self-supervised pretraining, tested on de-duplicated data with different numbers of labels, gave no gain over strong raw-feature baselines, and we report this as a negative result.

The Digital Twin works as an approval gate that protects the reward from poor proposals. A packet-level Simu5G sweep agrees with our calibrated QoS model on the order of delays but not on the shape of the delay curve. We also release a Digital Twin dataset of 393.6 million rows with six what-if scenarios, and an interactive monitor, Slice Twin, which is published on the web.

\vspace{0.4cm}
\noindent \textbf{Keywords:} network slicing, 6G, Digital Twin, self-supervised learning, temporal graph neural network, multi-agent reinforcement learning, traffic forecasting.
